\documentclass[../../../include/open-logic-section]{subfiles}

\begin{document}

\olfileid{sth}{ord-arithmetic}{expo}
\olsection{ক্রমসংখ্যার ঘাত}

এবার ক্রমসংখ্যার ঘাতের আলোচনা করি। দুর্ভাগ্যক্রমে
এর কোনো \emph{সহজ} সংশ্লেষমূলক সংজ্ঞা নেই।

অবশ্য সুস্পষ্ট সংশ্লেষমূলক সংজ্ঞা \emph{আছে}।
একটি উদাহরণ দিই। $\alpha \neq 0$ হলে
$\text{finfun}(\beta,\alpha)$ বলতে সেই সব অপেক্ষক
$f \colon \beta \to \alpha$-র সেট বোঝাই, যাদের জন্য
$\Setabs{\gamma \in \beta}{f(\gamma) \neq 0}$
কোনো স্বাভাবিক সংখ্যার সঙ্গে সমসংখ্যক।
$\text{finfun}(\beta,\alpha)$-তে $f \sqsubset g$
যদি এবং কেবল যদি $f \neq g$ এবং
$f(\gamma_0) < g(\gamma_0)$, যেখানে
$\gamma_0 = \text{max}\Setabs{\gamma \in \beta}{f(\gamma)
\neq g(\gamma)}$—এই শর্তে একটি সুক্রম সংজ্ঞায়িত করি।
তখন $\ordexpo{\alpha}{\beta}$-কে
$\ordtype{\text{finfun}(\beta, \alpha), \sqsubset}$
বলা যায়। পটার এই সুস্পষ্ট সংজ্ঞাটি ব্যবহার করেছেন,
তার পরেই ব্যাখ্যা করেছেন:
\begin{quote}
	এই ক্রমটি বেছে নেওয়ার একমাত্র কারণ, আমরা এমনভাবে
	ক্রমসংখ্যার ঘাত সংজ্ঞায়িত করতে চাই, যাতে যথাযথ
	পুনরাবৃত্তিমূলক শর্তটি মানা হয়\ldots; আর ক্রমযুক্ত
	যোগফল বা গুণফলের তুলনায় এটিকে কল্পনা করা অনেক কঠিন।
	\citep[পৃ.~১৯৯]{Potter2004}
\end{quote}
ঠিক তাই। যোগকে ব্যাখ্যা করেছি ``$\alpha$-র একটি
অনুলিপির পরে $\beta$-র একটি অনুলিপি'' হিসেবে; গুণকে
``$\alpha$-র অনুলিপিগুলির একটি $\beta$-ক্রম'' হিসেবে।
কিন্তু $\text{finfun}(\alpha, \gamma)$ নিয়ে তেমন
সংক্ষিপ্ত ব্যাখ্যা দেওয়া কঠিন। তাই ক্রমসংখ্যার ঘাতকে
সরাসরি \emph{ট্রান্সফাইনাইট পুনরাবৃত্তি দিয়ে}
সংজ্ঞায়িত করব, অর্থাৎ:

\begin{defn}\ollabel{ordexporecursion}
\begin{align*}
	\ordexpo{\alpha}{0} &= 1\\
	\ordexpo{\alpha}{\beta\ordplus 1} &=
	\ordexpo{\alpha}{\beta} \ordtimes \alpha\\
	\ordexpo{\alpha}{\beta} &=
	\bigcup_{0<\delta < \beta}\ordexpo{\alpha}{\delta}& &
	\text{যখন $\beta$ সীমা ক্রমসংখ্যা}
\end{align*}
% BN-SRC-805 TARGET-CORRECTION-BEGIN
\emph{উৎস-সংশোধন-নোট।} শূন্য ভিত্তিতে $0^0=1$,
কিন্তু প্রত্যেক ধনাত্মক ঘাতে $0^\beta=0$। সীমার
সংযোগে শূন্য ঘাতের মান রাখলে ভুল $0^\omega=1$ আসে।
তাই শুধু ধনাত্মক পূর্বঘাত নেওয়া হয়েছে। অশূন্য ভিত্তিতে
ঘাতগুলি একঘেয়ে এবং প্রথম ঘাত অন্তত $1$, তাই শূন্য ঘাত
বাদ দিলে সীমা-মান বদলায় না। ভিত্তি $1$-এ সব মান $1$;
এ কারণেও ঘাতে strict ঊর্ধ্বসীমা নেওয়া চলে না।
সংশ্লেষমূলক রূপটি শূন্য ভিত্তিতেও আলাদা করে প্রসারিত করা যায়:
খালি ক্ষেত্র থেকে একটিমাত্র খালি অপেক্ষক, অশূন্য ক্ষেত্র
থেকে শূন্য সেটে কোনও অপেক্ষক নেই; ক্রমপ্রকার যথাক্রমে $1$ ও $0$।
% BN-SRC-805 TARGET-CORRECTION-END
\end{defn}
% BN-SRC-451 TARGET-CORRECTION-BEGIN
\noindent\textbf{সম্পাদকীয় নোট।} মূল পাঠের
সংশ্লেষমূলক উদাহরণে $\text{finfun}(\alpha,\beta)$-এর
অপেক্ষকগুলির দিক উল্টো দেওয়া হয়েছে; $\beta \neq 0$
হলে সেটি $\ordexpo{\beta}{\alpha}$-র ক্রমপ্রকার দেয়। উপরে
$\alpha \neq 0$ শর্তে দিকটি সংশোধন করা হয়েছে।
উপরের সংশ্লেষমূলক উদাহরণ অশূন্য ভিত্তিতে বলা।
শূন্য ভিত্তির জন্য মূলের সীমা-পুনরাবৃত্তিও ঠিক করা
দরকার ছিল; সেই অতিরিক্ত মেরামত BN-SRC-805-এ স্পষ্ট।
% BN-SRC-451 TARGET-CORRECTION-END

সেটতাত্ত্বিক হিসেবেই কাজ করলে ক্রমসংখ্যার ঘাতের আরও
ধর্ম অনুসন্ধান করতে চাইতে পারি। কিন্তু আর বেশি কিছু
যোগ করার নেই, শুধু এই প্রত্যাশিত বিষয়টি লক্ষ করি:
ক্রমসংখ্যার ঘাত বিনিময়ী নয়। বস্তুত
$\ordexpo{2}{\omega} =
\bigcup_{\delta < \omega}\ordexpo{2}{\delta} = \omega$,
কিন্তু $\ordexpo{\omega}{2} = \omega \ordtimes \omega$।
ঘাত বিনিময়ী হবে এমন আশা করাও উচিত নয়, কারণ
স্বাভাবিক সংখ্যার ক্ষেত্রেও তা বিনিময়ী নয়:
$\ordexpo{2}{3} = 8 < 9 = \ordexpo{3}{2}$।

\begin{prob}
ট্রান্সফাইনাইট আবেশ ব্যবহার করে প্রমাণ করুন,
$\alpha \neq 0$ হলে
$\ordexpo{\alpha}{\beta} =
\ordtype{\text{finfun}(\beta, \alpha), \sqsubset}$
সংজ্ঞা নিলে \olref[sth][ord-arithmetic][expo]{ordexporecursion}-এর
পুনরাবৃত্তি-সমীকরণগুলি পাওয়া যায়।
\end{prob}

\end{document}
